\section{An irrational trace and failure of assembly}\label{sec:trace}

We now compute the trace of the projection constructed in
Proposition~\ref{prop:projection}.

\begin{proposition}\label{prop:irrational-trace}
The projection \(b_0\in C_r^*(\Gsur)\) satisfies
\begin{equation}\label{eq:irrational-trace}
\tau_{\Gsur}(b_0)=\sum_{i\ge1}p^{-k_i}\notin\Q.
\end{equation}
\end{proposition}

\begin{proof}
For a fixed \(\omega\), the diagonal entry of \(\pi_\omega(b_0)\) at
the identity is zero unless that vertex belongs to a valid copy. If it
does, the copy is unique and has some type \(i\); projection onto the
constants on that component has diagonal entry \(1/n_i\).

Exactly \(n_i\) indexed copies of type \(i\) contain the identity. Indeed,
their translating elements are precisely
\(g=\phi_i(v)^{-1}\) for the \(n_i\) distinct embedded vertices
\(v\in\operatorname{Vert}(\Theta_i)\).
Lemma~\ref{lem:distinct-copies} ensures that these indices give distinct
copies. By Lemma~\ref{lem:marginals}, each is valid on an event of Haar
measure \(p^{-k_i}\). Validity events for copies containing the identity
are mutually exclusive, including across types. The trace formula
\eqref{eq:trace-integral} and nonnegative countable summation therefore
give
\[
\tau_{\Gsur}(b_0)=\sum_{i\ge1}n_i\frac1{n_i}p^{-k_i}
=\sum_{i\ge1}p^{-k_i}.
\]
No independence of these events is needed.

Since \(p>2\) and the positive integers \(k_i\) strictly increase, this
number lies in \([0,1)\) and has base-\(p\) digits equal to one exactly
at positions \(k_i\), and zero elsewhere. There is no eventually
\((p-1)\) expansion ambiguity. The digits contain infinitely many ones,
but their successive gaps tend to infinity by
Proposition~\ref{prop:data}(iv). They cannot be eventually periodic.
Rational numbers have eventually periodic base-\(p\) expansions, so
\eqref{eq:irrational-trace} is irrational.
\end{proof}

\begin{proof}[Proof of Theorem~\ref{thm:main}]
The module construction gives a finitely generated discrete group
\(\Gsur\), and Proposition~\ref{prop:projection} gives an actual projection
\(b_0\in C_r^*(\Gsur)\). Its trace is positive and irrational by
Proposition~\ref{prop:irrational-trace}; in particular its integral
\(K_0\)-class is nonzero. Every class in the image of the proper-action
assembly map has rational trace by Theorem~\ref{thm:trace-image}. Thus
\([b_0]\) is outside that image.

The algebra \(C(\Omega)\rtimes_r\Gamma\) was used to represent
\(C_r^*(V\rtimes\Gamma)\) and to analyze the operator. The assembly map
just obstructed is the coefficient-free reduced map for \(\Gsur\) itself,
with the full proper classifying space. Its failure of surjectivity in
degree zero refutes the universal Baum--Connes assertion.
\end{proof}

The obstruction persists after rationalization. Extend the canonical
trace to \(K_0(C_r^*(\Gsur))\otimes_{\Z}\Q\) by
\(x\otimes q\mapsto q\tau_{\Gsur*}(x)\). Every element in the image
of \(\mu_0^{\Gsur}\otimes_{\Z}\operatorname{id}_{\Q}\) has rational
trace, since it is a finite rational linear combination of assembly
classes. The trace of \([b_0]\otimes1\) remains irrational, so this class
is outside the image of rationalized assembly as well.
